\chapter{Test et résultats}

\section{Introduction}
Ce chapitre présente les tests réalisés pour valider la fonctionnalité 
de recueil de feedbacks et de réclamations développée dans le cadre 
du stage chez Bewize, ainsi que les résultats obtenus. La phase de 
test constitue une étape indispensable dans tout cycle de développement 
logiciel, car elle permet de s'assurer que la solution développée 
répond correctement aux besoins exprimés et fonctionne de manière 
fiable dans toutes les situations prévues.

Les tests ont été réalisés en collaboration avec le Tech Lead et 
le Product Manager de Bewize, en suivant une approche structurée 
couvrant plusieurs niveaux de test : les tests unitaires, les tests 
d'intégration et les tests fonctionnels bout en bout. Cette approche 
multi-niveaux garantit une couverture complète de la solution et 
une confiance maximale dans la qualité du livrable final.

%==========================================================
\section{Stratégie de test}
%==========================================================

\subsection{Niveaux de test adoptés}
La stratégie de test adoptée pour ce projet s'appuie sur la pyramide 
de tests, un modèle reconnu dans l'industrie du développement logiciel 
qui préconise une répartition équilibrée entre les différents types 
de tests. Cette pyramide recommande d'avoir un grand nombre de tests 
unitaires à la base, un nombre intermédiaire de tests d'intégration 
au milieu, et un nombre réduit de tests fonctionnels bout en bout 
au sommet.

\begin{figure}[H]
    \centering
    %\includegraphics[width=0.75\textwidth]{Figures/pyramide_tests.png}
    \caption{Pyramide de tests adoptée pour le projet}
    \label{fig:pyramide_tests}
\end{figure}

Les trois niveaux de test mis en œuvre sont les suivants :

\begin{itemize}
    \item \textbf{Tests unitaires :} ils testent chaque composant 
    ou fonction de manière isolée, sans dépendance avec les autres 
    composants du système. Ces tests sont rapides à exécuter et 
    faciles à maintenir.
    
    \item \textbf{Tests d'intégration :} ils vérifient la bonne 
    communication et interaction entre les différents composants 
    du système, notamment entre le front-end et le back-end, et 
    entre le back-end et la base de données.
    
    \item \textbf{Tests fonctionnels bout en bout :} ils simulent 
    le comportement réel d'un utilisateur en testant le parcours 
    complet depuis l'application mobile jusqu'au dashboard monitor, 
    en passant par la WebView et l'API.
\end{itemize}

\subsection{Environnement de test}
Les tests ont été réalisés dans un environnement dédié, distinct 
de l'environnement de développement et de production. Cet environnement 
de test utilise une base de données PostgreSQL séparée, peuplée avec 
des données de test représentatives des cas d'utilisation réels.

\begin{table}[H]
\centering
\renewcommand{\arraystretch}{1.5}
\begin{tabular}{|p{4cm}|p{8cm}|}
\hline
\textbf{Composant} & \textbf{Configuration de test} \\
\hline
Base de données & PostgreSQL de test avec données fictives \\
\hline
Back-end Spring Boot & Instance de test avec profil "test" activé \\
\hline
WebView React JS & Build de test accessible via localhost \\
\hline
Dashboard monitor & Build de test accessible via localhost \\
\hline
Application mobile & Émulateur Android et simulateur iOS \\
\hline
Outil de test API & Postman avec collection de tests dédiée \\
\hline
\end{tabular}
\caption{Environnement de test}
\label{tab:env_test}
\end{table}

%==========================================================
\section{Tests unitaires}
%==========================================================

\subsection{Tests unitaires du back-end Spring Boot}
Les tests unitaires du back-end ont été développés en utilisant 
JUnit 5, le framework de test standard pour les applications Java, 
combiné avec Mockito pour la simulation des dépendances. Ces tests 
couvrent principalement la couche service, qui contient la logique 
métier de l'application.

Les méthodes testées incluent :
\begin{itemize}
    \item \textbf{Création d'un feedback valide :} vérification que 
    le statut est initialisé à NOUVEAU, que la date d'envoi est 
    correctement générée et que les données sont sauvegardées
    \item \textbf{Création avec données invalides :} vérification 
    que les exceptions appropriées sont levées en cas de données 
    manquantes ou incorrectes
    \item \textbf{Récupération d'un feedback par ID :} vérification 
    du retour correct ou de l'exception ResourceNotFoundException
    \item \textbf{Mise à jour du statut :} vérification que les 
    transitions de statut valides sont acceptées et les invalides rejetées
    \item \textbf{Calcul des statistiques :} vérification de l'exactitude 
    des comptages par type et par statut
\end{itemize}

\begin{table}[H]
\centering
\renewcommand{\arraystretch}{1.8}
\begin{tabular}{|c|p{5cm}|p{2.5cm}|p{2.5cm}|}
\hline
\textbf{ID} & \textbf{Test} & \textbf{Résultat attendu} & 
\textbf{Résultat} \\
\hline
UT01 & Création feedback avec données valides & 
Feedback créé avec statut NOUVEAU & Succès \\
\hline
UT02 & Création feedback avec message vide & 
Exception levée & Succès \\
\hline
UT03 & Création feedback avec type null & 
Exception levée & Succès \\
\hline
UT04 & Récupération feedback existant par ID & 
Feedback retourné & Succès \\
\hline
UT05 & Récupération feedback inexistant & 
ResourceNotFoundException & Succès \\
\hline
UT06 & Mise à jour statut NOUVEAU vers EN\_COURS & 
Statut mis à jour & Succès \\
\hline
UT07 & Mise à jour statut EN\_COURS vers TRAITE & 
Statut mis à jour & Succès \\
\hline
UT08 & Calcul statistiques avec données variées & 
Comptages exacts & Succès \\
\hline
\end{tabular}
\caption{Résultats des tests unitaires back-end}
\label{tab:tests_unitaires}
\end{table}

\begin{figure}[H]
    \centering
    %\includegraphics[width=0.9\textwidth]{Figures/junit_results.png}
    \caption{Résultats des tests unitaires JUnit}
    \label{fig:junit_results}
\end{figure}

\subsection{Tests unitaires du front-end React JS}
Les tests unitaires du front-end ont été développés en utilisant 
Jest et React Testing Library, les outils de test standards pour 
les applications React JS. Ces tests vérifient le comportement 
correct de chaque composant React de manière isolée.

Les composants testés incluent :
\begin{itemize}
    \item \textbf{Composant formulaire :} vérification du rendu 
    correct, de la gestion des changements de valeurs et de la 
    soumission du formulaire
    \item \textbf{Validation des champs :} vérification que les 
    messages d'erreur appropriés sont affichés en cas de données invalides
    \item \textbf{Message de confirmation :} vérification que le 
    message de succès s'affiche correctement après l'envoi
    \item \textbf{Gestion des erreurs réseau :} vérification que 
    le message d'erreur s'affiche correctement en cas d'échec de l'API
\end{itemize}

%==========================================================
\section{Tests d'intégration}
%==========================================================

\subsection{Tests d'intégration de l'API REST}
Les tests d'intégration de l'API REST ont été réalisés avec Postman, 
un outil professionnel de test d'APIs. Une collection Postman complète 
a été créée pour tester l'ensemble des endpoints de manière systématique 
et reproductible.

\begin{figure}[H]
    \centering
    %\includegraphics[width=0.9\textwidth]{Figures/postman_collection.png}
    \caption{Collection Postman des tests d'intégration}
    \label{fig:postman_collection}
\end{figure}

\subsubsection{Test POST /api/feedbacks}
Ce test vérifie la création correcte d'un nouveau feedback ou 
réclamation via l'API REST. Plusieurs scénarios ont été testés :

\begin{table}[H]
\centering
\renewcommand{\arraystretch}{1.8}
\begin{tabular}{|c|p{4cm}|p{2.5cm}|p{2.5cm}|p{1.5cm}|}
\hline
\textbf{ID} & \textbf{Scénario} & \textbf{Données envoyées} & 
\textbf{Résultat attendu} & \textbf{Statut} \\
\hline
IT01 & Feedback valide & Type + Message & HTTP 201 & Succès \\
\hline
IT02 & Réclamation valide & Type + Message & HTTP 201 & Succès \\
\hline
IT03 & Message vide & Type uniquement & HTTP 400 & Succès \\
\hline
IT04 & Type manquant & Message uniquement & HTTP 400 & Succès \\
\hline
IT05 & Corps de requête vide & Aucune donnée & HTTP 400 & Succès \\
\hline
\end{tabular}
\caption{Tests d'intégration — POST /api/feedbacks}
\label{tab:test_post}
\end{table}

\subsubsection{Test GET /api/feedbacks}
Ce test vérifie la récupération correcte de la liste des feedbacks 
et réclamations, avec et sans filtres appliqués :

\begin{table}[H]
\centering
\renewcommand{\arraystretch}{1.8}
\begin{tabular}{|c|p{4.5cm}|p{3cm}|p{2cm}|}
\hline
\textbf{ID} & \textbf{Scénario} & \textbf{Résultat attendu} & 
\textbf{Statut} \\
\hline
IT06 & Liste sans filtre & Tous les retours & Succès \\
\hline
IT07 & Filtre par type FEEDBACK & Feedbacks uniquement & Succès \\
\hline
IT08 & Filtre par type RECLAMATION & Réclamations uniquement & Succès \\
\hline
IT09 & Filtre par statut NOUVEAU & Retours NOUVEAU & Succès \\
\hline
IT10 & Filtre par statut TRAITE & Retours TRAITE & Succès \\
\hline
\end{tabular}
\caption{Tests d'intégration — GET /api/feedbacks}
\label{tab:test_get}
\end{table}

\subsubsection{Test PUT /api/feedbacks/\{id\}/statut}
Ce test vérifie la mise à jour correcte du statut d'un feedback :

\begin{table}[H]
\centering
\renewcommand{\arraystretch}{1.8}
\begin{tabular}{|c|p{4.5cm}|p{3cm}|p{2cm}|}
\hline
\textbf{ID} & \textbf{Scénario} & \textbf{Résultat attendu} & 
\textbf{Statut} \\
\hline
IT11 & Passage à EN\_COURS & Statut mis à jour & Succès \\
\hline
IT12 & Passage à TRAITE & Statut mis à jour & Succès \\
\hline
IT13 & ID inexistant & HTTP 404 & Succès \\
\hline
IT14 & Statut invalide & HTTP 400 & Succès \\
\hline
\end{tabular}
\caption{Tests d'intégration — PUT statut}
\label{tab:test_put}
\end{table}

\subsection{Tests d'intégration base de données}
Les tests d'intégration de la base de données ont vérifié la 
persistance correcte des données dans PostgreSQL :

\begin{itemize}
    \item Vérification que les feedbacks sont correctement enregistrés 
    avec tous leurs attributs dans la table feedback
    \item Vérification de l'intégrité des données après mise à jour 
    du statut de traitement
    \item Vérification des contraintes de la base de données (NOT NULL, 
    types de données)
    \item Vérification des performances des requêtes de filtrage 
    avec un volume de données représentatif
\end{itemize}

\begin{figure}[H]
    \centering
    %\includegraphics[width=0.88\textwidth]{Figures/pgadmin_data.png}
    \caption{Données de test dans PostgreSQL via pgAdmin}
    \label{fig:pgadmin_data}
\end{figure}

%==========================================================
\section{Tests fonctionnels bout en bout}
%==========================================================

\subsection{Présentation des tests fonctionnels}
Les tests fonctionnels bout en bout simulent le comportement réel 
des utilisateurs et valident le parcours complet de la fonctionnalité. 
Ces tests ont été réalisés en collaboration avec le Tech Lead et 
validés par le Product Manager de Bewize lors de la Sprint Review 
du dernier Sprint.

\subsection{Scénario 1 — Envoi d'un feedback depuis l'application mobile}

\textbf{Objectif :} Vérifier que l'utilisateur peut envoyer un 
feedback depuis l'application mobile et que les données apparaissent 
correctement dans le dashboard monitor.

\textbf{Étapes du test :}
\begin{enumerate}
    \item Lancer l'application mobile Bewize sur un émulateur Android
    \item Localiser et cliquer sur le bouton feedback dans l'interface
    \item Vérifier que la WebView s'ouvre correctement
    \item Sélectionner le type "Feedback"
    \item Saisir un message de test
    \item Cliquer sur le bouton Envoyer
    \item Vérifier que le message de confirmation s'affiche
    \item Accéder au dashboard monitor
    \item Vérifier que le feedback apparaît dans la liste avec le 
    statut NOUVEAU
\end{enumerate}

\textbf{Résultat :} $\checkmark$ Test réussi — Le feedback est correctement 
envoyé et affiché dans le dashboard avec toutes ses informations.

\begin{figure}[H]
    \centering
    %\includegraphics[width=0.85\textwidth]{Figures/test_feedback.png}
    \caption{Test fonctionnel — Envoi d'un feedback}
    \label{fig:test_feedback}
\end{figure}

\subsection{Scénario 2 — Envoi d'une réclamation et traitement}

\textbf{Objectif :} Vérifier le parcours complet d'une réclamation, 
depuis son envoi par l'utilisateur mobile jusqu'à son traitement 
complet par l'équipe Bewize dans le dashboard monitor.

\textbf{Étapes du test :}
\begin{enumerate}
    \item Envoyer une réclamation depuis la WebView mobile
    \item Vérifier la réception dans le dashboard monitor
    \item Changer le statut de NOUVEAU à EN\_COURS
    \item Vérifier la mise à jour dans le dashboard
    \item Changer le statut de EN\_COURS à TRAITE
    \item Vérifier le statut final dans la liste
\end{enumerate}

\textbf{Résultat :} $\checkmark$ Test réussi — La réclamation est traitée 
correctement avec les transitions de statut fonctionnelles.

\begin{figure}[H]
    \centering
    %\includegraphics[width=0.85\textwidth]{Figures/test_reclamation.png}
    \caption{Test fonctionnel — Traitement d'une réclamation}
    \label{fig:test_reclamation}
\end{figure}

\subsection{Scénario 3 — Validation du formulaire}

\textbf{Objectif :} Vérifier que la validation des champs fonctionne 
correctement et empêche l'envoi de données invalides.

\textbf{Étapes du test :}
\begin{enumerate}
    \item Ouvrir le formulaire WebView
    \item Cliquer sur Envoyer sans remplir aucun champ
    \item Vérifier que les messages de validation s'affichent
    \item Remplir uniquement le type sans le message
    \item Vérifier que le message de validation s'affiche
    \item Remplir tous les champs correctement
    \item Vérifier que l'envoi réussit
\end{enumerate}

\textbf{Résultat :} $\checkmark$ Test réussi — La validation fonctionne 
correctement dans tous les cas testés.

\subsection{Scénario 4 — Filtrage dans le dashboard}

\textbf{Objectif :} Vérifier que les filtres du dashboard monitor 
fonctionnent correctement et retournent les résultats attendus.

\textbf{Étapes du test :}
\begin{enumerate}
    \item Envoyer plusieurs feedbacks et réclamations de test
    \item Accéder au dashboard monitor
    \item Appliquer le filtre "Type : Feedback"
    \item Vérifier que seuls les feedbacks s'affichent
    \item Appliquer le filtre "Statut : NOUVEAU"
    \item Vérifier que seuls les retours NOUVEAU s'affichent
    \item Combiner les deux filtres
    \item Vérifier que les résultats correspondent aux deux critères
\end{enumerate}

\textbf{Résultat :} $\checkmark$ Test réussi — Les filtres fonctionnent 
correctement individuellement et en combinaison.

\begin{figure}[H]
    \centering
    %\includegraphics[width=0.85\textwidth]{Figures/test_filtre.png}
    \caption{Test fonctionnel — Filtrage dans le dashboard}
    \label{fig:test_filtre}
\end{figure}

\subsection{Tableau récapitulatif des tests fonctionnels}

\begin{table}[H]
\centering
\renewcommand{\arraystretch}{1.8}
\begin{tabular}{|c|p{5cm}|p{2.5cm}|p{2.5cm}|}
\hline
\textbf{ID} & \textbf{Scénario testé} & \textbf{Résultat attendu} & 
\textbf{Résultat obtenu} \\
\hline
TF01 & Envoi feedback depuis mobile & Données en base, confirmation & 
Conforme \\
\hline
TF02 & Envoi réclamation depuis mobile & Données en base, confirmation & 
Conforme \\
\hline
TF03 & Validation champs vides & Messages d'erreur affichés & 
Conforme \\
\hline
TF04 & Consultation dashboard & Liste complète affichée & 
Conforme \\
\hline
TF05 & Filtre par type & Résultats filtrés correctement & 
Conforme \\
\hline
TF06 & Filtre par statut & Résultats filtrés correctement & 
Conforme \\
\hline
TF07 & Changement statut NOUVEAU à EN\_COURS & Statut mis à jour & 
Conforme \\
\hline
TF08 & Changement statut EN\_COURS à TRAITE & Statut mis à jour & 
Conforme \\
\hline
TF09 & Affichage statistiques & Données calculées correctement & 
Conforme \\
\hline
TF10 & Test sur émulateur Android & Fonctionnel sur Android & 
Conforme \\
\hline
TF11 & Test sur simulateur iOS & Fonctionnel sur iOS & 
Conforme \\
\hline
\end{tabular}
\caption{Tableau récapitulatif des tests fonctionnels}
\label{tab:tests_fonctionnels}
\end{table}

%==========================================================
\section{Résultats et bilan}
%==========================================================

\subsection{Résultats globaux}
L'ensemble des tests réalisés ont abouti à des résultats très 
satisfaisants, confirmant la qualité et la fiabilité de la 
solution développée. Le taux de réussite global des tests 
est de 100\%, avec tous les scénarios testés conformes 
aux résultats attendus.

\begin{table}[H]
\centering
\renewcommand{\arraystretch}{1.8}
\begin{tabular}{|p{4cm}|c|c|c|}
\hline
\textbf{Type de test} & \textbf{Tests réalisés} & 
\textbf{Tests réussis} & \textbf{Taux de réussite} \\
\hline
Tests unitaires back-end & 8 & 8 & 100\% \\
\hline
Tests unitaires front-end & 6 & 6 & 100\% \\
\hline
Tests d'intégration API & 14 & 14 & 100\% \\
\hline
Tests fonctionnels & 11 & 11 & 100\% \\
\hline
\textbf{Total} & \textbf{39} & \textbf{39} & \textbf{100\%} \\
\hline
\end{tabular}
\caption{Bilan global des tests}
\label{tab:bilan_tests}
\end{table}

\begin{figure}[H]
    \centering
    %\includegraphics[width=0.8\textwidth]{Figures/bilan_tests.png}
    \caption{Représentation graphique du bilan des tests}
    \label{fig:bilan_tests}
\end{figure}

\subsection{Résultats attendus atteints}
À l'issue des tests, tous les résultats attendus définis dans le 
cahier des charges ont été atteints :

\begin{itemize}
    \item \textbf{Fonctionnalité de recueil de feedbacks :} $\checkmark$ 
    Les utilisateurs peuvent envoyer facilement leurs feedbacks 
    depuis l'application mobile via la WebView
    
    \item \textbf{Fonctionnalité de recueil de réclamations :} $\checkmark$ 
    Les utilisateurs peuvent envoyer leurs réclamations en quelques 
    clics depuis l'application mobile
    
    \item \textbf{Page WebView accessible depuis le mobile :} $\checkmark$ 
    La WebView s'ouvre correctement dans l'application mobile sans 
    redirection vers un navigateur externe
    
    \item \textbf{Transmission des données vers le back-end :} $\checkmark$ 
    Les données sont transmises de manière fiable et sécurisée via 
    les APIs REST Spring Boot
    
    \item \textbf{Affichage dans le dashboard monitor :} $\checkmark$ 
    L'équipe Bewize peut consulter, filtrer et traiter les retours 
    reçus depuis le dashboard monitor interne
    
    \item \textbf{Parcours testé de bout en bout :} $\checkmark$ 
    Le parcours complet a été validé avec succès depuis l'envoi 
    par l'utilisateur jusqu'au traitement par l'équipe Bewize
\end{itemize}

\subsection{Validation par le Product Manager}
La solution développée a été présentée au Product Manager de Bewize 
lors de la Sprint Review finale. Cette démonstration a couvert 
l'ensemble du parcours utilisateur : l'envoi d'un feedback depuis 
l'application mobile, la réception dans le dashboard monitor et 
le traitement de la réclamation. Le Product Manager a validé 
la solution et confirmé qu'elle répond pleinement aux attentes 
et aux besoins exprimés dans le cahier des charges.

\begin{figure}[H]
    \centering
    %\includegraphics[width=0.85\textwidth]{Figures/validation_pm.png}
    \caption{Démonstration de la solution au Product Manager}
    \label{fig:validation_pm}
\end{figure}

%==========================================================
\section{Conclusion}
%==========================================================

Ce chapitre a présenté l'ensemble des tests réalisés pour valider 
la fonctionnalité de recueil de feedbacks et de réclamations 
développée chez Bewize, ainsi que les résultats obtenus.

La stratégie de test adoptée, couvrant trois niveaux complémentaires 
(tests unitaires, tests d'intégration et tests fonctionnels bout 
en bout), a permis de valider la solution de manière exhaustive 
et rigoureuse. Le taux de réussite de 100\% sur l'ensemble des 
39 tests réalisés confirme la qualité et la fiabilité de 
l'implémentation.

L'ensemble des résultats attendus définis dans le cahier des charges 
ont été atteints et validés par le Product Manager de Bewize lors 
de la Sprint Review finale. La solution est fonctionnelle, fiable 
et prête à être déployée en production dans l'environnement Bewize.